\documentclass[xcolor={x11names}]{beamer}
\usetheme{Madrid}

\usepackage{amssymb}
\usepackage{ulem}
\usepackage[utf8]{inputenc}
\usepackage{mathtools}
\usepackage{multicol}
%\usepackage[x11names]{xcolor}
\usefonttheme{professionalfonts}


% Subfigures
\usepackage{caption}
\usepackage{subcaption}

% Check/cross mark
\usepackage{pifont}% http://ctan.org/pkg/pifont
\newcommand{\cmark}{\ding{51}}%
\newcommand{\xmark}{\ding{55}}%


% License
\usepackage[
    type={CC},
    modifier={by-nc-sa},
    version={4.0},
    imagewidth=4pt
]{doclicense}




% Change base colour beamer@blendedblue (originally RGB: 0.2,0.2,0.7)
\colorlet{beamer@blendedblue}{DarkSeaGreen4}







%% MATH commands
\DeclareMathOperator{\Var}{Var}


%% THEOREMS
%\newtheorem{theorem}{Theorem}
\newtheorem{thm}{Teorema}[section] % the main one
% Definición
%\theoremstyle{definition}
\newtheorem{definicion}{Definición}[section]
\newtheorem{lema}{Lema}[section]



%% PFGplots %%
\usepackage{pgfplots}

%% Exponential distribution
\pgfmathdeclarefunction{exponential}{1}{%
  \pgfmathparse{(#1)*exp(-#1*x)}%
}
\pgfmathdeclarefunction{exponentialcdf}{1}{%
  \pgfmathparse{1-exp(-#1*x)}%
}

%% Poisson distribution
\pgfmathdeclarefunction{poiss}{1}{%
  \pgfmathparse{(#1^x)*exp(-#1)/(x!)}%
}

%% Normal distribution (#1=mu, #2=sigma)
% John D. Cook approx. https://tex.stackexchange.com/a/124629
\pgfmathdeclarefunction{normalcdf}{2}{%
  \pgfmathparse{1/(1 + exp(-0.07056*((x-#1)/#2)^3 - 1.5976*(x-#1)/#2))}%
}




\newcommand{\red}[1]{{\color{red}#1}}
\newcommand{\blue}[1]{{\color{blue}#1}}

%%%%%%%%%%
%% TIKZ %%
%%%%%%%%%%
\usepackage{tikz}
\usepackage{animate}
\usetikzlibrary{positioning}
\usetikzlibrary{shapes,arrows, positioning, calc}
\usetikzlibrary{overlay-beamer-styles}
\usetikzlibrary{chains,shapes.multipart}
\usetikzlibrary{scopes}
\usetikzlibrary{automata}
\usetikzlibrary{positioning}  %                 ...positioning nodes
\usetikzlibrary{arrows}       %                 ...customizing arrows
\usetikzlibrary{intersections}


%%%%%%%%%
%% PGF %%
%%%%%%%%%
\usepgfplotslibrary{fillbetween}


%%% Insert section name before the section %%%
\AtBeginSection[]{
  \begin{frame}
  \vfill
  \centering
  \begin{beamercolorbox}[sep=8pt,center,shadow=true,rounded=true]{title}
    \usebeamerfont{title}\insertsectionhead\par%
  \end{beamercolorbox}
  \vfill
  \end{frame}
}



\title[Tema 5]{Tema 5: Teletráfico en Redes de Datos}
%% \subtitle{Redes y Servicios de Telecomunicaciones (RSTC)\\
%% Grado en Ingeniería de Tecnologías y Servicios de Telecomunicación}
%\author{M. Saiful Bari\inst{1} \and Mr X\inst{2}}

\titlegraphic{%
\includegraphics[height=2.5cm]{../tema4/figs/RSTC-grande.png}\\%
\doclicenseIcon {\tiny \hspace{1em}\doclicenseText}\\%
\href{https://github.com/MartinPJorge/RSTC-queueing-slides}{\includegraphics[height=1cm]{../tema4/figs/github-logo.png}}%
}


\author{\textcolor{white}{RSTC curso 2023-2024}}
%\author{Jorge Martín Pérez\inst{1}}
%\institute{
%    \inst{1}
%    Departamento de Ingeniería Telemática, Universidad Politécnica de Madrid
%}

\date{\today}







%%%%%%%%%%%%%%%%%%%%
%%% SLIDES START %%%
%%%%%%%%%%%%%%%%%%%%
\begin{document}


%%% TITLE %%%
\frame{\titlepage}


\begin{frame}{Contenido}
    \tableofcontents
\end{frame}




\section{Introducción}
\begin{frame}{\secname}
    Hemos visto colas M/M/1
    \begin{figure}
        
\begin{tikzpicture}[start chain=going right,>=latex,node distance=0pt]
% the rectangular shape with vertical lines
\node[rectangle split, rectangle split parts=6,
draw, rectangle split horizontal,text height=1cm,text depth=0.5cm,on chain,inner ysep=0pt] (wa) {};
\fill[white] ([xshift=-\pgflinewidth,yshift=-\pgflinewidth]wa.north west) rectangle ([xshift=-15pt,yshift=\pgflinewidth]wa.south);

% the circle
\node[draw,circle,on chain,minimum size=1.5cm] (se) {$\mu$};

% the arrows and labels
\draw[->] (se.east) -- +(20pt,0);
\draw[<-] (wa.west) -- +(-20pt,0) node[left] {$\lambda$};
\node[align=center,below] at (wa.south) {Waiting \\ Area};
\node[align=center,below] at (se.south) {Service \\ Node};
\end{tikzpicture}

    \end{figure}
    con tiempos:
    \begin{itemize}
        \item de llegada exponenciales
            $t_l\sim Exp(\lambda)$
        \item de servicio exponenciales
            $t_s\sim Exp(\mu)$
    \end{itemize}

    \pause
    \vfill

    {\color{blue}
    Pero, ¿y si el tiempo de servicio $t_s$
    sigue otra distribución?
    }
     
    \pause
    \begin{itemize}
        {\color{red}
        \item sistema M/G/1
        }
    \end{itemize}
\end{frame}



\begin{frame}{\secname}
    Hemos estudiado una sola cola
    \begin{figure}
        
\begin{tikzpicture}[start chain=going right,>=latex,node distance=0pt]
% the rectangular shape with vertical lines
\node[rectangle split, rectangle split parts=6,
draw, rectangle split horizontal,text height=1cm,text depth=0.5cm,on chain,inner ysep=0pt] (wa) {};
\fill[white] ([xshift=-\pgflinewidth,yshift=-\pgflinewidth]wa.north west) rectangle ([xshift=-15pt,yshift=\pgflinewidth]wa.south);

% the circle
\node[draw,circle,on chain,minimum size=1.5cm] (se) {$\mu$};

% the arrows and labels
\draw[->] (se.east) -- +(20pt,0);
\draw[<-] (wa.west) -- +(-20pt,0) node[left] {$\lambda$};
\node[align=center,below] at (wa.south) {Waiting \\ Area};
\node[align=center,below] at (se.south) {Service \\ Node};
\end{tikzpicture}

    \end{figure}

    \vfill
    {\color{blue}
    Pero, ¿y si hay más colas?
    }
    \pause
    \begin{itemize}
        {\color{red}
        \item redes de Jackson
        }
    \end{itemize}
    \begin{figure}
        \begin{tikzpicture}[color=Firebrick2,thick]


% the rectangular shape with vertical lines
\node[rectangle split, rectangle split parts=6,
draw, rectangle split horizontal,text height=1cm,text depth=0.5cm,inner ysep=0pt] (wa) {};
\fill[white] ([xshift=-\pgflinewidth,yshift=-\pgflinewidth]wa.north west) rectangle ([xshift=-15pt,yshift=\pgflinewidth]wa.south);

% the circle
\node[draw,circle,minimum size=1.5cm,anchor=west] (se) at (wa.east) {$\mu$};

% the arrows and labels
\draw[->] (se.east) -- +(20pt,0) node (seEnd) {};
\draw[<-] (wa.west) -- +(-20pt,0) node[left] {$\lambda$};




%%%%%%%%%%%%%%%%%%%
%%% Second queue
%%%%%%%%%%%%%%%%%%%

% the rectangular shape with vertical lines
\node[rectangle split, rectangle split parts=6,
    draw, rectangle split horizontal,text height=1cm,text depth=0.5cm,inner ysep=0pt,anchor=west] (wa2) at (seEnd.east) {};
\fill[white] ([xshift=-\pgflinewidth,yshift=-\pgflinewidth]wa2.north west) rectangle ([xshift=-15pt,yshift=\pgflinewidth]wa2.south);

% the circle
\node[draw,circle,minimum size=1.5cm,anchor=west] (se2) at (wa2.east) {$\mu$};

% the arrows and labels
\draw[->] (se2.east) -- +(20pt,0) node (seEnd2) {};


\end{tikzpicture}

    \end{figure}
\end{frame}


\begin{frame}{Contenido}
    \tableofcontents
\end{frame}


\section{Sistema M/G/1}
\subsection{No Markoviano}
\begin{frame}{\secname: \subsecname}
    Tiempo de servicio sigue una distribución
    general\footnote{Por ejemplo,
    $G(\mu)=U(\tfrac{1}{2\mu}, \tfrac{2}{3\mu})$}
    $t_s\sim G(\mu)$.

    \begin{figure}
        
\begin{tikzpicture}[start chain=going right,>=latex,node distance=0pt]
% the rectangular shape with vertical lines
\node[rectangle split, rectangle split parts=6,
draw, rectangle split horizontal,text height=1cm,text depth=0.5cm,on chain,inner ysep=0pt] (wa) {};
\fill[white] ([xshift=-\pgflinewidth,yshift=-\pgflinewidth]wa.north west) rectangle ([xshift=-15pt,yshift=\pgflinewidth]wa.south);

% the circle
\node[draw,circle,on chain,minimum size=1.5cm] (se) {$\mu$};

% the arrows and labels
\draw[->] (se.east) -- +(20pt,0);
\draw[<-] (wa.west) -- +(-20pt,0) node[left] {$\lambda$};
\node[align=center,below] at (wa.south) {Waiting \\ Area};
\node[align=center,below] at (se.south) {Service \\ Node};
\end{tikzpicture}

    \end{figure}

    \vfill
    Para modelar como cadena de Markov
    es necesario que
    \begin{itemize}
        \item tiempo estancia en estado
            $t_i\sim Exp(\nu_i)$.
    \end{itemize}
\end{frame}



\begin{frame}{\secname: \subsecname}
    \begin{figure}
        %%% \tikzset{node distance=1cm, % Minimum distance between two nodes. Change if necessary.
%%%          every state/.style={ % Sets the properties for each state
%%%            semithick,draw=HotPink3!50,
%%%            fill=HotPink3!20},
%%%          initial text={},     % No label on start arrow
%%%          double distance=4pt, % Adjust appearance of accept states
%%%          every edge/.style={  % Sets the properties for each transition
%%%          draw, ->,>=stealth',     % Makes edges directed with bold arrowheads
%%%            auto, thick},
%%% }
\begin{tikzpicture}
   

\draw[|->] (0,0) -- (5,0) node[anchor=west] {$t$};

% Estados
\node[state,fill=HotPink3!20] (current)
    at (0,1.3) {$i=3$};
\node[dashed,state,fill=DodgerBlue4!20] (arrive)
    at (2,2) {$i=4$};
\node[dashed,state,fill=DodgerBlue1!20] (exit)
    at (4,2) {$i=2$};

% Posibles eventos
\draw[thick,dashed,<-] (2,0) to (arrive.south);
\draw[thick,dashed,->] (4,0) to (exit.south);

% Tiempos hasta eventos
\draw[|-|,thick,gray,DodgerBlue4] (0,0.2) -- node[pos=.5,fill=white,inner sep=0] {$t_l$} (1.95,0.2);
\draw[|-|,thick,gray,DodgerBlue1] (0,0.5) -- node[pos=.8,fill=white,inner sep=0] {$t_s$} (3.95,0.5);

\node (legend) at (3,3) 
    {\tiny posibles transiciones};
\draw[->] (legend.east)
    to[out=0,in=90] (exit.north);
\draw[->] (legend.west) 
    to[out=180,in=90] (arrive.north);
    
\end{tikzpicture}

    \end{figure}

    \vfill
    Veamos si se cumple que $t_i\sim Exp(\nu_i)$:
    \begin{multline}
        \mathbb{P}(t_i>\tau) =
        \mathbb{P}(\min\{t_l,t_s\}>\tau)=
        \mathbb{P}(t_l>\tau)
        \mathbb{P}(t_s>\tau)\\
        =
        \left(1 - \frac{\tau}{\mu} \right)
        e^{-\lambda\tau} \neq e^{-\nu_i \tau}
    \end{multline}
    con $t_s\sim G(\mu)=U\left(\tfrac{1}{2\mu},
    \tfrac{2}{3\mu}\right),\
    \tau\in\left[\tfrac{1}{2\mu},
    \tfrac{2}{3\mu}\right]$.
\end{frame}




\subsection{Tiempo medio de espera en cola}
\begin{frame}{\secname: \subsecname}
    Podemos obtener el tiempo medio
    de espera en cola $\mathbb{E}[W(t)]$ de
    un M/G/1.
    \begin{figure}
        
\begin{tikzpicture}[start chain=going right,>=latex,node distance=0pt]
% the rectangular shape with vertical lines
\node[rectangle split, rectangle split parts=6,
draw, rectangle split horizontal,text height=1cm,text depth=0.5cm,on chain,inner ysep=0pt] (wa) {};
\fill[white] ([xshift=-\pgflinewidth,yshift=-\pgflinewidth]wa.north west) rectangle ([xshift=-15pt,yshift=\pgflinewidth]wa.south);

% the circle
\node[draw,circle,on chain,minimum size=1.5cm] (se) {$\mu$};

% the arrows and labels
\draw[->] (se.east) -- +(20pt,0);
\draw[<-] (wa.west) -- +(-20pt,0) node[left] {$\lambda$};
\node[align=center,below] at (wa.south) {Waiting \\ Area};
\node[align=center,below] at (se.south) {Service \\ Node};
\end{tikzpicture}

    \end{figure}
    Veamos lo que espera un usuario
    \emph{nuevo}:
    \begin{enumerate}
        \item $\mathbb{E}[Q(t)]\tfrac{1}{\mu}$
            en cola ; y
        \item la media del tiempo residual
            $R$ del que
            se está sirviendo.
    \end{enumerate}
    \begin{align*}
        \mathbb{E}[W(t)]&=
        \mathbb{E}[Q(t)]\tfrac{1}{\mu}
        + \mathbb{E}[R(t)]\\
        &\implies \mathbb{E}[(W(t))]=
        \frac{\mathbb{E}[R(t)]}{1-\rho}
    \end{align*}
\end{frame}




\begin{frame}{\secname: \subsecname}
    Interpretación gráfica de tiempo
    residual en cada instante $R(t)$:
    \begin{figure}
        \begin{tikzpicture}
    \draw[->] (0,0) -- (5,0) node[anchor=west]
        {$t$};
    \draw[->] (0,0) -- (0,2.5)
        node[anchor=south] {$R(t)$};

    % yticks
    \draw (-.1,1) node[anchor=east]
        {1} -- (.1,1);
    \draw (-.1,2) node[anchor=east]
        {2} -- (.1,2);

    % xticks
    \draw (1,-.1) node[anchor=north]
        {1} -- (1,.1);
    \draw (2,-.1) node[anchor=north]
        {2} -- (2,.1);
    \draw (3,-.1) node[anchor=north]
        {3} -- (3,.1);
    \draw (4,-.1) node[anchor=north]
        {4} -- (4,.1);
    
    % Areas residuales
    \path [fill=Firebrick1!20,draw=black]
        (3,0) -- (3,1) -- (4,0);
    \path[fill=DodgerBlue1!20,draw=black]
        (1,0) -- (1,2) -- (3,0);
    \node[align=center] at (2,.3) {$u_1$};
    \node[align=center] at (3.3,.3) {$u_2$};




    % Llegadas
    \draw[|->] (0,-2) -- (5,-2) node[anchor=west]
        {$t$};
    \draw[->,color=DodgerBlue1,very thick]
        (1,-1.5) node[anchor=south] {$u_1$}
        -- (1,-2);
    \draw[->,very thick,color=Firebrick1]
        (1.5,-1.5) node[anchor=south] (u2)
        {$u_2$} -- (1.5,-2);
    \draw (1.5,-2) -- (1.5,-2.2)
        node[anchor=north,color=Firebrick1,
        align=left]
        {$t_2=1.5$};

    % Residual for second user
    \node[circle,fill=Firebrick1,
       minimum size=4,inner sep=0,draw=black]
       (dotResidual) at (1.5,1.5) {};
    \node[anchor=west] at (dotResidual.east)
       {$R(t_2)=\frac{3}{2}$};

   % Link second user residual
   \draw[dashed,color=Firebrick1]
       (u2.north) -- (dotResidual);


\end{tikzpicture}


    \end{figure}
\end{frame}



\begin{frame}{\secname: \subsecname}
    \footnote{
    Ilustración y demostración basadas
    en~\cite[Figura 8.3]{amable}.}
    La media del tiempo residual corresponde
    con el promedio de áreas.
    \begin{figure}
        \begin{tikzpicture}
    \draw (0,0) -- (7,0)
        node[anchor=west] {$\ldots$};
    \draw[->] (7.75,0) -- (9.5,0)
        node[anchor=west] {$t$};
    \draw[->] (0,0) -- (0,2.5)
        node[anchor=south] {$R(t)$};

    % yticks
    \draw (-.1,1) node[anchor=east]
        {1} -- (.1,1);
    \draw (-.1,2) node[anchor=east]
        {2} -- (.1,2);

    % xticks
    \foreach \i in {1,...,6} {
        \draw (\i,-.1) node[anchor=north]
            {\i} -- (\i,.1);
    }
    
    % Areas residuales
    \path [fill=Firebrick1!20,draw=black]
        (3,0) -- (3,1) -- (4,0);
    \path[fill=DodgerBlue1!20,draw=black]
        (1,0) -- (1,2) -- (3,0);
    \path[fill=Gold1!20,draw=black]
        (4.5,0) -- (4.5,.5) -- (5,0);
    \path[fill=SpringGreen4!20,draw=black]
        (5,0) -- (5,1.5) -- (6.5,0);
    \path[fill=Orchid1!20,draw=black]
        (8,0) -- (8,.75) -- (8.75,0);
    \node[align=center] at (2,.3) {$u_1$};
    \node[align=center] at (3.3,.3) {$u_2$};
    \node[align=center] at (4.7,.125)
        {\tiny $u_3$};
    \node[align=center] at (5.5,.3) {$u_4$};
    \node[align=center] at (8.3,.2)
        {$u_\infty$};


    \node[circle,minimum size=3,inner sep=0,
        fill=DodgerBlue1,draw=black] (timeU1)
        at (1,2) {};
    \node[anchor=west] at (timeU1.east)
        {$t_1=2$};



\end{tikzpicture}


    \end{figure}
    \begin{align*}
        \mathbb{E}[R(t)]=&\lim_{t\to\infty}
        \frac{1}{t}\int_0^t
        R(\tau)\ d\tau\\
        =& \lim_{t\to\infty} \frac{1}{t}
        \sum_{i=1}^{N(t)} \frac{t_{s,i}^2}{2}
    \end{align*}
\end{frame}




\begin{frame}{\secname: \subsecname}
    Multiplicando y dividiendo por
    $N(t)$:
    \begin{align*}
        \mathbb{E}[R(t)]
        =& \lim_{t\to\infty}\frac{1}{t}
        \sum_{i=1}^{N(t)} \frac{t_{s,i}^2}{2}\\
        =& \frac{1}{2}
        \left( \lim_{t\to\infty} \frac{N(t)}{t}
        \right)
        \left( \lim_{t\to\infty} \frac{1}{N(t)}
        \sum_{i=1}^{N(t)} t_{s,i}^2
        \right)\\
        =& \frac{1}{2} \cdot \lambda\cdot
        \mathbb{E}[t_s^2]
    \end{align*}
    con $t_{s,i}$ la realización
    de la v.a. del tiempo de servicio
    para el usuario $i$.
\end{frame}



\begin{frame}{\secname: \subsecname}
    \begin{lema}[Fórmula de Pollaczek-Khintchine]
        El tiempo medio de espera en cola
        de en un sistema M/G/1 es
        \begin{equation}
            \mathbb{E}[W(t)]=
            \frac{\lambda\mathbb{E}[t_s^2]}{2(1-\rho)}
        \end{equation}
        con $t_s$ la v.a. del tiempo de
        servicio que sigue una distribución
        $G$.
    \end{lema}
    \textit{Demostración}:
    \begin{equation*}
        \mathbb{E}[W(t)]
        =\frac{\mathbb{E}[R(t)]}{1-\rho}
        =\frac{\tfrac{1}{2}\lambda\mathbb{E}[t_s^2]}{1-\rho}
        =\frac{\lambda\mathbb{E}[t_s^2]}{2(1-\rho)}
    \end{equation*}

\end{frame}




\subsection{Ejemplos distribuciones servicio}
\begin{frame}{\secname: \subsecname}
    \textit{Ejemplo}: supongamos un tiempo
    de servicio $t_s\sim Exp(\mu)$.
    \begin{align*}
        \mathbb{E}[t_s^2]
        =& \int_0^\infty
        \tau^2 \mu e^{-\mu\tau}
        \ d\tau\\
        \underbrace{=}_{partes}&
        \left[-\tau^2e^{-\mu\tau} \right]_{\tau=0}^\infty
        -\int_0^\infty -e^{-\mu\tau}2\tau\ 
        d\tau\\
        =&\int_0^\infty
        2\tau e^{-\mu\tau}\ d\tau\\
        \underbrace{=}_{partes}&
        \frac{2}{\mu^2}
    \end{align*}
    Usando Pollaczek-Khintchine tenemos
    \begin{equation*}
        \mathbb{E}[W(t)]=
        \frac{\lambda\mathbb{E}[t_s^2]}{2(1-\rho)}
        =\frac{\rho}{\mu(1-\rho)}
    \end{equation*}
    la expresión que vimos para M/M/1.
\end{frame}






\begin{frame}{\secname: \subsecname}
    \textit{Ejemplo}: supongamos un tiempo
    de servicio $t_s\sim U(0,\frac{2}{\mu})$.
    \begin{align*}
        \mathbb{E}[t_s^2]
        =& \int_0^{\tfrac{2}{\mu}}
        \tau^2 \frac{1}{2/\mu}
        \ d\tau\\
        =& \frac{\mu}{2}\left[
            \frac{\tau^3}{3}
        \right]_{\tau=0}^{\tfrac{2}{\mu}}\\
        =& \frac{4}{3\mu^2}
    \end{align*}
    Usando Pollaczek-Khintchine tenemos
    \begin{equation*}
        \mathbb{E}[W(t)]=
        \frac{\lambda\mathbb{E}[t_s^2]}{2(1-\rho)}
        =\frac{2}{3}\frac{\rho}{\mu(1-\rho)}
    \end{equation*}
\end{frame}



\begin{frame}{\secname: \subsecname}
    Otra manera de ver el momento de
    segundo orden es sabiendo que
    \begin{equation*}
        \mathbb{E}[t_s^2]=\Var[t_s]+
        \mathbb{E}^2[t_s]
    \end{equation*}
    \vfill
    \textit{Ejemplo}:
    \begin{itemize}
        \item $t_s\sim Exp(\mu)\implies
            \mathbb{E}[t_s^2]=\frac{1}{\mu^2}+
            \frac{1}{\mu^2}$;
        \item $t_s\sim U\left(0,\tfrac{2}{\mu}
            \right)
            \implies
            \mathbb{E}[t_s^2]=\frac{1}{12}
            \left(\frac{2}{\mu}-0\right)^2+
            \frac{1}{\mu^2}
            =\frac{4}{3\mu^2}$.
    \end{itemize}
    que coincide con las expresiones anteriores.
\end{frame}


\subsection{M/M/1 como caso peor}
\begin{frame}{\secname: \subsecname}
    El tiempo medio de servicio
    $\mathbb{E}[T(t)]$
    es pesimista en un M/M/1.

    \vfill


    \begin{figure}
        \begin{tikzpicture}

\begin{axis}[every axis plot post/.append style={
  mark=none,domain=0:.9,samples=40},
  axis x line*=bottom,
  axis y line*=left,
  ymin=0,
  ymax=8,
  xmax=1,
  xlabel={$\lambda$ [usuarios/sec]},
  ylabel={$\mathbb{E}[T(t)]$},
  enlargelimits=upper,
  width=4in,
  height=2in]
  \addplot[ultra thick,smooth,color=HotPink2] {2/3*x/(1-x)+1} node [pos=.25,anchor=north west]
      {M/U/1};
  \addplot[ultra thick,smooth,color=HotPink4] {x/(1-x)+1} node [pos=.2,anchor=south east]
      {M/M/1};
  \addplot +[mark=none,dashed] coordinates {(1,0) (1,8)};
\end{axis}


\end{tikzpicture}


    \end{figure}
    \textit{Ejemplo (arriba)}: el tiempo medio
    total es menor\footnote{
        Tomamos $\mu=1$
    [usuario/sec].} en una
    uniforme.


    % \begin{align*}
    %     \mathbb{E}[W_U(t)]+\tfrac{1}{\mu}<&
    %     \mathbb{E}[W_{Exp}(t)]+\tfrac{1}\mu\\
    %     \frac{2}{3}\frac{\rho}{\mu(1-\rho)}<&
    %     \frac{\rho}{\mu(1-\rho)}\\
    %     \frac{2}{3}<&1
    % \end{align*}
    % vemos que es
\end{frame}






%   Sabiendo que:
%   \begin{itemize}
%       \item hay $Q(t)$ usuarios en
%           cola; y
%       \item tiempo medio de servicio
%           $\mathbb{E}[t_s]=\tfrac{1}{\mu}$
%   \end{itemize}

%   Un usuario nuevo tarda en ser servido
%   $Q(t)\cdot \tfrac{1}{\mu} + 







\section{Redes de Colas}
\begin{frame}{\secname}
    Podemos estudiar cómo modelar una red
    de colas (e.g., routers).
    \vfill
    \begin{figure}
        \usetikzlibrary{calc}
\begin{tikzpicture}


% the rectangular shape with vertical lines
\node[rectangle split, rectangle split parts=6,
draw, rectangle split horizontal,text height=1cm,text depth=0.5cm,inner ysep=0pt] (wa) {};
\fill[white] ([xshift=-\pgflinewidth,yshift=-\pgflinewidth]wa.north west) rectangle ([xshift=-15pt,yshift=\pgflinewidth]wa.south);

% the circle
\node[draw,circle,minimum size=1.5cm,anchor=west] (se) at (wa.east) {$\mu_1$};

% the arrows and labels
\draw[->] (se.east) -- +(20pt,0) node (seEnd) {};
\draw[<-] (wa.west) -- +(-20pt,0) node[left] {$\lambda_1$};




%%%%%%%%%%%%%%%%%%%
%%% Second queue
%%%%%%%%%%%%%%%%%%%

% the rectangular shape with vertical lines
\node[rectangle split, rectangle split parts=6,
    draw, rectangle split horizontal,text height=1cm,text depth=0.5cm,inner ysep=0pt,anchor=west] (wa2) at (seEnd.east) {};
\fill[white] ([xshift=-\pgflinewidth,yshift=-\pgflinewidth]wa2.north west) rectangle ([xshift=-15pt,yshift=\pgflinewidth]wa2.south);

% the circle
\node[draw,circle,minimum size=1.5cm,anchor=west] (se2) at (wa2.east) {$\mu_2$};

% the arrows and labels
\draw[->] (se2.east) -- +(40pt,0) node (seEnd2) {};



%%%%%%%%%%%%%%%%%%%
%%% Third queue
%%%%%%%%%%%%%%%%%%%

% the rectangular shape with vertical lines
\node[rectangle split, rectangle split parts=6,
    draw, rectangle split horizontal,text height=1cm,text depth=0.5cm,inner ysep=0pt,anchor=west] (wa3) at ($(wa2.south)-(2.5,1.5)$) {};
\fill[white]
    ([xshift=+\pgflinewidth,yshift=-\pgflinewidth]wa3.north east)
    rectangle ([xshift=-1pt,yshift=\pgflinewidth]wa3.south);
% the circle
\node[draw,circle,minimum size=1.5cm,anchor=east] (se3) at (wa3.west) {$\mu_3$};



%%%%%%%%%%%%%%%%%
%%% DRAW JOINS
%%%%%%%%%%%%%%%%%
\node[circle,inner sep=0pt,minimum size=3pt,draw=black,fill=black] (joinOut) at ($(se2.east)+(10pt,0)$) {};
\node[circle,inner sep=0pt,minimum size=3pt,draw=black,fill=black] (joinIn) at ($(wa.west)-(10pt,0)$) {};
\node[circle,inner sep=0pt,minimum size=3pt,draw=black,fill=black] (joinMid) at ($(wa2.west)-(15pt,0)$) {};
\node[circle,inner sep=0pt,minimum size=3pt,draw=black,fill=black] (joinMidDown) at ($(wa3.east)+(20pt,0)$) {};


% the exit
\draw[->] (joinOut) |- node[pos=0.3,anchor=west]
    {$p_{2,3}$} (wa3.east);

% the midfeedback
\draw[->] (joinMid) --
    ($(joinMid) - (0,1.1)$) -|
    node[pos=0.8,anchor=west] {$p_{1,3}$}
    (joinMidDown.north);

% the out feedback
\draw[->] (se3.west) -| (joinIn.south);



\end{tikzpicture}

    \end{figure}
    con $p_{2,3}$ la probabilidad de ir
    de la cola 2 a la 3.
\end{frame}


\begin{frame}{\secname}
    Las llegadas $\lambda_i$ a la cola
    $i$ se averiguan usando probabilidades
    $p_{i,j}$.

    \vfill

    \textit{Ejemplo}:
    \begin{equation*}
        {\color{HotPink4}\lambda_3} = p_{1,3}\lambda_{1,o}
        + p_{2,3}\lambda_{2,o}
    \end{equation*}
    con $\lambda_{i,o}$ la tasa de salidas
    de la cola $i$.
    \begin{figure}
        \resizebox{.7\textwidth}{!}{%
            \usetikzlibrary{calc}
\begin{tikzpicture}


% the rectangular shape with vertical lines
\node[rectangle split, rectangle split parts=6,
draw, rectangle split horizontal,text height=1cm,text depth=0.5cm,inner ysep=0pt] (wa) {};
\fill[white] ([xshift=-\pgflinewidth,yshift=-\pgflinewidth]wa.north west) rectangle ([xshift=-15pt,yshift=\pgflinewidth]wa.south);

% the circle
\node[draw,circle,minimum size=1.5cm,anchor=west] (se) at (wa.east) {$\mu_1$};

% the arrows and labels
\draw[->] (se.east) -- +(20pt,0) node (seEnd) {};
\draw[<-] (wa.west) -- +(-20pt,0) node[left] {};




%%%%%%%%%%%%%%%%%%%
%%% Second queue
%%%%%%%%%%%%%%%%%%%

% the rectangular shape with vertical lines
\node[rectangle split, rectangle split parts=6,
    draw, rectangle split horizontal,text height=1cm,text depth=0.5cm,inner ysep=0pt,anchor=west] (wa2) at (seEnd.east) {};
\fill[white] ([xshift=-\pgflinewidth,yshift=-\pgflinewidth]wa2.north west) rectangle ([xshift=-15pt,yshift=\pgflinewidth]wa2.south);

% the circle
\node[draw,circle,minimum size=1.5cm,anchor=west] (se2) at (wa2.east) {$\mu_2$};

% the arrows and labels
\draw[->] (se2.east) -- +(40pt,0)
    node[pos=.7,anchor=south west ](seEnd2)
    {$(1-p_{2,3})$};



%%%%%%%%%%%%%%%%%%%
%%% Third queue
%%%%%%%%%%%%%%%%%%%

% the rectangular shape with vertical lines
\node[rectangle split, rectangle split parts=6,
    draw, rectangle split horizontal,text height=1cm,text depth=0.5cm,inner ysep=0pt,anchor=west] (wa3) at ($(wa2.south)-(2.5,1.5)$) {};
\fill[white]
    ([xshift=+\pgflinewidth,yshift=-\pgflinewidth]wa3.north east)
    rectangle ([xshift=-1pt,yshift=\pgflinewidth]wa3.south);
% the circle
\node[draw,circle,minimum size=1.5cm,anchor=east] (se3) at (wa3.west) {$\mu_3$};



%%%%%%%%%%%%%%%%%
%%% DRAW JOINS
%%%%%%%%%%%%%%%%%
\node[circle,inner sep=0pt,minimum size=3pt,draw=black,fill=black] (joinOut) at ($(se2.east)+(10pt,0)$) {};
\node[circle,inner sep=0pt,minimum size=3pt,draw=black,fill=black] (joinIn) at ($(wa.west)-(10pt,0)$) {};
\node[circle,inner sep=0pt,minimum size=3pt,draw=black,fill=black] (joinMid) at ($(wa2.west)-(15pt,0)$) {};
\node[circle,inner sep=0pt,minimum size=3pt,draw=black,fill=black] (joinMidDown) at ($(wa3.east)+(20pt,0)$) {};


% entrance lambda1
\node[anchor=south west] at (joinIn.east) {$\lambda_1$};


% the exit
\draw[->] (joinOut) |- node[pos=0.3,anchor=west]
    {$p_{2,3}$} (wa3.east);

% the midfeedback
\draw[->] (joinMid) --
    ($(joinMid) - (0,1.1)$) -|
    node[pos=0.8,anchor=west] {$p_{1,3}$}
    (joinMidDown.north);

% the out feedback
\draw[->] (se3.west) -| (joinIn.south);



%%%%%%%%%%%%%%%%%
%%% ARRIVALS
%%%%%%%%%%%%%%%%%
\node[anchor=west,circle,
    draw=HotPink4, inner sep=1,
    fill=HotPink4!20,] at
    (wa2.west)
    {$\lambda_2$};
\node[anchor=east,circle,
    draw=HotPink4, inner sep=1,
    fill=HotPink4!20,] at
    ($(wa3.east)-(2pt,0)$)
    {$\lambda_3$};


\end{tikzpicture}
%
        }
    \end{figure}
    
\end{frame}


\subsection{Distribución salida M/M/1}
\begin{frame}{\secname: \subsecname}
    \begin{lema}[Tiempo entre salidas]
        En un sistema M/M/1 el tiempo entre
        salidas y llegadas siguen la misma
        distribución. Es decir,
        se cumple:
        \begin{align*}
            t_e\sim&  Exp(\lambda)\\
            t_l\sim&  Exp(\lambda)
        \end{align*}
    \end{lema}
    \textit{Demostración}:
    \begin{multline*}
        F_{t_e}(\tau)=
        \pi_0 \mathbb{P}(t_l+t_s\leq\tau)+
        (1-\pi_0) \mathbb{P}(t_s\leq\tau)\\
        =(1-\rho) \int_0^\tau
        f_{t_l+t_s}(t)\ dt
        + \rho (1-e^{-\mu\tau})\\
        =(1-\rho)\left(
        1-\frac{\mu}{\mu-\lambda}
        e^{-\lambda\tau}
        +\frac{\lambda}{\mu-\lambda}
        e^{-\mu\tau}
        \right)
        + \rho (1-e^{-\mu\tau})\\
        = 1-e^{-\lambda\tau}
    \end{multline*}
\end{frame}



\begin{frame}{\secname: \subsecname}
    \textit{Ejemplo (cont.)}: sabiendo
    que $\lambda_{i,o}=\lambda_i$
    en la red de colas de abajo tendríamos que
    \begin{align*}
        {\color{HotPink4}\lambda_3}=&
        p_{1,3}\lambda_1+p_{2,3}\lambda_2\\
        =&p_{1,3}\lambda_1+p_{2,3}
        (1-p_{1,3})\lambda_1
    \end{align*}

    \begin{figure}
        \resizebox{.7\textwidth}{!}{%
            \usetikzlibrary{calc}
\begin{tikzpicture}


% the rectangular shape with vertical lines
\node[rectangle split, rectangle split parts=6,
draw, rectangle split horizontal,text height=1cm,text depth=0.5cm,inner ysep=0pt] (wa) {};
\fill[white] ([xshift=-\pgflinewidth,yshift=-\pgflinewidth]wa.north west) rectangle ([xshift=-15pt,yshift=\pgflinewidth]wa.south);

% the circle
\node[draw,circle,minimum size=1.5cm,anchor=west] (se) at (wa.east) {$\mu_1$};

% the arrows and labels
\draw[->] (se.east) -- +(20pt,0) node (seEnd) {};
\draw[<-] (wa.west) -- +(-20pt,0) node[left] {};




%%%%%%%%%%%%%%%%%%%
%%% Second queue
%%%%%%%%%%%%%%%%%%%

% the rectangular shape with vertical lines
\node[rectangle split, rectangle split parts=6,
    draw, rectangle split horizontal,text height=1cm,text depth=0.5cm,inner ysep=0pt,anchor=west] (wa2) at (seEnd.east) {};
\fill[white] ([xshift=-\pgflinewidth,yshift=-\pgflinewidth]wa2.north west) rectangle ([xshift=-15pt,yshift=\pgflinewidth]wa2.south);

% the circle
\node[draw,circle,minimum size=1.5cm,anchor=west] (se2) at (wa2.east) {$\mu_2$};

% the arrows and labels
\draw[->] (se2.east) -- +(40pt,0)
    node[pos=.7,anchor=south west ](seEnd2)
    {$(1-p_{2,3})$};



%%%%%%%%%%%%%%%%%%%
%%% Third queue
%%%%%%%%%%%%%%%%%%%

% the rectangular shape with vertical lines
\node[rectangle split, rectangle split parts=6,
    draw, rectangle split horizontal,text height=1cm,text depth=0.5cm,inner ysep=0pt,anchor=west] (wa3) at ($(wa2.south)-(2.5,1.5)$) {};
\fill[white]
    ([xshift=+\pgflinewidth,yshift=-\pgflinewidth]wa3.north east)
    rectangle ([xshift=-1pt,yshift=\pgflinewidth]wa3.south);
% the circle
\node[draw,circle,minimum size=1.5cm,anchor=east] (se3) at (wa3.west) {$\mu_3$};



%%%%%%%%%%%%%%%%%
%%% DRAW JOINS
%%%%%%%%%%%%%%%%%
\node[circle,inner sep=0pt,minimum size=3pt,draw=black,fill=black] (joinOut) at ($(se2.east)+(10pt,0)$) {};
\node[circle,inner sep=0pt,minimum size=3pt,draw=black,fill=black] (joinIn) at ($(wa.west)-(10pt,0)$) {};
\node[circle,inner sep=0pt,minimum size=3pt,draw=black,fill=black] (joinMid) at ($(wa2.west)-(15pt,0)$) {};
\node[circle,inner sep=0pt,minimum size=3pt,draw=black,fill=black] (joinMidDown) at ($(wa3.east)+(20pt,0)$) {};


% entrance lambda1
\node[anchor=south west] at (joinIn.east) {$\lambda_1$};


% the exit
\draw[->] (joinOut) |- node[pos=0.3,anchor=west]
    {$p_{2,3}$} (wa3.east);

% the midfeedback
\draw[->] (joinMid) --
    ($(joinMid) - (0,1.1)$) -|
    node[pos=0.8,anchor=west] {$p_{1,3}$}
    (joinMidDown.north);

% the out feedback
\draw[->] (se3.west) -| (joinIn.south);



%%%%%%%%%%%%%%%%%
%%% ARRIVALS
%%%%%%%%%%%%%%%%%
\node[anchor=west,circle,
    draw=HotPink4, inner sep=1,
    fill=HotPink4!20,] at
    (wa2.west)
    {$\lambda_2$};
\node[anchor=east,circle,
    draw=HotPink4, inner sep=1,
    fill=HotPink4!20,] at
    ($(wa3.east)-(2pt,0)$)
    {$\lambda_3$};


\end{tikzpicture}
%
        }
    \end{figure}
\end{frame}



\subsection{Vector de Probabilidades}
\begin{frame}{\secname: \subsecname}
    Quiero saber la probabilidad de tener
    $\mathbf{K}=(3,1,4)$ usuarios en
    equilibrio:
    \begin{equation*}
        \mathbb{P}(\mathbf{K})=
        \left(\mathbb{P}(N_1=3),
        \mathbb{P}(N_2=1),
        \mathbb{P}(N_3=4)\right)
    \end{equation*}
    con $N_i$ el número de usuarios en la
    cola $i$ al alcanzar equilibrio.

    \vfill

    \begin{figure}
        \resizebox{.7\textwidth}{!}{%
            \usetikzlibrary{calc}
\begin{tikzpicture}


% the rectangular shape with vertical lines
\node[rectangle split, rectangle split parts=6,
draw, rectangle split horizontal,
rectangle split part fill={white,white,white,white,HotPink3!20,HotPink3!20},
text height=1cm,text depth=0.5cm,inner ysep=0pt] (wa) {};
\fill[white] ([xshift=-\pgflinewidth,yshift=-\pgflinewidth]wa.north west) rectangle ([xshift=-15pt,yshift=\pgflinewidth]wa.south);

% the circle
\node[draw,circle,minimum size=1.5cm,anchor=west,fill=HotPink3!20] (se) at (wa.east) {$\mu_1$};

% the arrows and labels
\draw[->] (se.east) -- +(20pt,0) node (seEnd) {};
\draw[<-] (wa.west) -- +(-20pt,0) node[left] {};




%%%%%%%%%%%%%%%%%%%
%%% Second queue
%%%%%%%%%%%%%%%%%%%

% the rectangular shape with vertical lines
\node[rectangle split, rectangle split parts=6,
    draw, rectangle split horizontal,text height=1cm,text depth=0.5cm,inner ysep=0pt,anchor=west] (wa2) at (seEnd.east) {};
\fill[white] ([xshift=-\pgflinewidth,yshift=-\pgflinewidth]wa2.north west) rectangle ([xshift=-15pt,yshift=\pgflinewidth]wa2.south);

% the circle
\node[draw,circle,minimum size=1.5cm,anchor=west,fill=HotPink3!20] (se2) at (wa2.east) {$\mu_2$};

% the arrows and labels
\draw[->] (se2.east) -- +(40pt,0) node (seEnd2) {};



%%%%%%%%%%%%%%%%%%%
%%% Third queue
%%%%%%%%%%%%%%%%%%%

% the rectangular shape with vertical lines
\node[rectangle split, rectangle split parts=6,
    draw, rectangle split horizontal,
    rectangle split part fill={HotPink3!20,
    HotPink3!20,HotPink3!20,white},
    text height=1cm,text depth=0.5cm,inner ysep=0pt,anchor=west] (wa3) at ($(wa2.south)-(2.5,1.5)$) {};
\fill[white]
    ([xshift=+\pgflinewidth,yshift=-\pgflinewidth]wa3.north east)
    rectangle ([xshift=-1pt,yshift=\pgflinewidth]wa3.south);
% the circle
\node[draw,circle,minimum size=1.5cm,anchor=east,fill=HotPink3!20] (se3) at (wa3.west) {$\mu_3$};



%%%%%%%%%%%%%%%%%
%%% DRAW JOINS
%%%%%%%%%%%%%%%%%
\node[circle,inner sep=0pt,minimum size=3pt,draw=black,fill=black] (joinOut) at ($(se2.east)+(10pt,0)$) {};
\node[circle,inner sep=0pt,minimum size=3pt,draw=black,fill=black] (joinIn) at ($(wa.west)-(10pt,0)$) {};
\node[circle,inner sep=0pt,minimum size=3pt,draw=black,fill=black] (joinMid) at ($(wa2.west)-(15pt,0)$) {};
\node[circle,inner sep=0pt,minimum size=3pt,draw=black,fill=black] (joinMidDown) at ($(wa3.east)+(20pt,0)$) {};


% entrance lambda1
\node[anchor=south west] at (joinIn.east) {$\lambda_1$};


% the exit
\draw[->] (joinOut) |- node[pos=0.3,anchor=west]
    {$p_{2,3}$} (wa3.east);

% the midfeedback
\draw[->] (joinMid) --
    ($(joinMid) - (0,1.1)$) -|
    node[pos=0.8,anchor=west] {$p_{1,3}$}
    (joinMidDown.north);

% the out feedback
\draw[->] (se3.west) -| (joinIn.south);



\end{tikzpicture}

        }
    \end{figure}
\end{frame}




\subsection{Redes de Jackson}
\begin{frame}{\secname: \subsecname}
    Para que exista $\mathbb{P}(\mathbf{K})$
    necesitamos que sea una red de Jackson:

    \vfill

    \begin{definicion}[Red de Jackson]
        Las colas interconectadas forman una
        red de Jackson si:
        \begin{enumerate}
            \item la red es abierta, y las
                llegadas en cola
                siguen un proceso de Poisson;
            \item los tiempos de servicio son
                exponenciales y las colas
                son FIFO;
            \item un usuario pasa de la
                cola $i$ a $j$ con probabilidad
                $p_{i,j}$, o sale con
                probabilidad
                $1-\sum_{j}p_{i,j}$; y
            \item la carga de las colas está
                por debajo de 1, i.e.
                $\rho_i<1,\ \forall i$
        \end{enumerate}
    \end{definicion}
\end{frame}




\begin{frame}{\secname: \subsecname}
    \textit{Ejemplo}: ¿tenemos una red
    de Jackson en esta diapositiva?

    \begin{figure}
        \resizebox{.7\textwidth}{!}{%
            \usetikzlibrary{calc}
\begin{tikzpicture}


% the rectangular shape with vertical lines
\node[rectangle split, rectangle split parts=6,
draw, rectangle split horizontal,text height=1cm,text depth=0.5cm,inner ysep=0pt] (wa) {};
\fill[white] ([xshift=-\pgflinewidth,yshift=-\pgflinewidth]wa.north west) rectangle ([xshift=-15pt,yshift=\pgflinewidth]wa.south);

% the circle
\node[draw,circle,minimum size=1.5cm,anchor=west] (se) at (wa.east) {$\mu_1$};

% the arrows and labels
\draw[->] (se.east) -- +(20pt,0) node (seEnd) {};
\draw[<-] (wa.west) -- +(-20pt,0) node[left] {$\lambda_1$};




%%%%%%%%%%%%%%%%%%%
%%% Second queue
%%%%%%%%%%%%%%%%%%%

% the rectangular shape with vertical lines
\node[rectangle split, rectangle split parts=6,
    draw, rectangle split horizontal,text height=1cm,text depth=0.5cm,inner ysep=0pt,anchor=west] (wa2) at (seEnd.east) {};
\fill[white] ([xshift=-\pgflinewidth,yshift=-\pgflinewidth]wa2.north west) rectangle ([xshift=-15pt,yshift=\pgflinewidth]wa2.south);

% the circle
\node[draw,circle,minimum size=1.5cm,anchor=west] (se2) at (wa2.east) {$\mu_2$};

% the arrows and labels
\draw[->] (se2.east) -- +(40pt,0) node (seEnd2) {};



%%%%%%%%%%%%%%%%%%%
%%% Third queue
%%%%%%%%%%%%%%%%%%%

% the rectangular shape with vertical lines
\node[rectangle split, rectangle split parts=6,
    draw, rectangle split horizontal,text height=1cm,text depth=0.5cm,inner ysep=0pt,anchor=west] (wa3) at ($(wa2.south)-(2.5,1.5)$) {};
\fill[white]
    ([xshift=+\pgflinewidth,yshift=-\pgflinewidth]wa3.north east)
    rectangle ([xshift=-1pt,yshift=\pgflinewidth]wa3.south);
% the circle
\node[draw,circle,minimum size=1.5cm,anchor=east] (se3) at (wa3.west) {$\mu_3$};



%%%%%%%%%%%%%%%%%
%%% DRAW JOINS
%%%%%%%%%%%%%%%%%
\node[circle,inner sep=0pt,minimum size=3pt,draw=black,fill=black] (joinOut) at ($(se2.east)+(10pt,0)$) {};
\node[circle,inner sep=0pt,minimum size=3pt,draw=black,fill=black] (joinIn) at ($(wa.west)-(10pt,0)$) {};
\node[circle,inner sep=0pt,minimum size=3pt,draw=black,fill=black] (joinMid) at ($(wa2.west)-(15pt,0)$) {};
\node[circle,inner sep=0pt,minimum size=3pt,draw=black,fill=black] (joinMidDown) at ($(wa3.east)+(20pt,0)$) {};


% the exit
\draw[->] (joinOut) |- node[pos=0.3,anchor=west]
    {$p_{2,3}$} (wa3.east);

% the midfeedback
\draw[->] (joinMid) --
    ($(joinMid) - (0,1.1)$) -|
    node[pos=0.8,anchor=west] {$p_{1,3}$}
    (joinMidDown.north);

% the out feedback
\draw[->] (se3.west) -| (joinIn.south);



\end{tikzpicture}

        }
    \end{figure}

    Comprobemos las propiedades:
    \begin{enumerate}
        \item red abierta y
            ¿llegadas de Poisson?

            \begin{itemize}
                \item \cmark
                    cola $i=1$ con llegadas
                    $Poiss(\lambda_1)$;
                \item \textbf{?}
                    cola $i=2$ recibe
                    $(1-p_{1,3})$ de salidas
                    de cola $i=1$; y
                \item \textbf{?}
                    cola $i=3$ recibe
                    $p_{1,3}$ de salidas
                    de cola $i=1$, y el
                    $p_{2,3}$ de salidas de
                    cola $i=2$;
            \end{itemize}
    \end{enumerate}
\end{frame}




\begin{frame}{\secname: \subsecname}
    \begin{lema}[Descomposición proceso
        de Poisson]
        Si $\{N(t), t>0\}$ es un proceso
        estocástico que sigue una distribución
        $Poiss(\lambda)$,
        y cada llegada se asigna a
        \begin{itemize}
            \item $N_1(t)$ con probabilidad $p$,
                y
            \item $N_2(t)$ con probabilidad
                $1-p$;
        \end{itemize}
        entonces:
        \begin{itemize}
            \item $N_1(t)\sim Poiss(p\lambda)$
            \item $N_2(t)\sim Poiss\left((1-p)\lambda\right)$
        \end{itemize}
    \end{lema}

    \footnote{
        Figura basada en
        \cite[Figura 3.7]{amable}} 
    \begin{figure}    
        \begin{tikzpicture}
    \draw[|->] (0,0) -- (10,0) node[anchor=west] {time [min]};

    % Time axis top
    \draw[|->] (0,0) -- (10,0) node[anchor=west] {time [min]};

    % ticks
    \foreach \t in {1,...,9}
        \draw (\t,.05) -- (\t,-.05) node[pos=1,anchor=north] {$\t$};

    
    % arrivals A
    \foreach \t in {2,4} {
        \draw[->,color=Firebrick1,very thick] (\t,.5) node[anchor=west,rotate=70] {compacto}  --(\t,0);
        \draw[->,color=Firebrick1,very thick] (\t,-1.5) node[anchor=west,rotate=70] {}  --(\t,-2);
    }

    % arrivals B
    \foreach \t in {.1,3,3.5,5,5.5,7,8,9.2} {
        \draw[->,color=Firebrick4,very thick] (\t,.5) node[anchor=west,rotate=70] {}  --(\t,0);
        \draw[->,color=Firebrick4,very thick] (\t,-2.5) node[anchor=west,rotate=70] {}  --(\t,-3);
    }




    % Time axis A
    \draw[|->] (0,-2) -- (10,-2) node[anchor=west] {time [min]};

    % Time axis B
    \draw[|->] (0,-3) -- (10,-3) node[anchor=west] {time [min]};
\end{tikzpicture}

    \end{figure}
\end{frame}


\begin{frame}{\secname: \subsecname}
    \textit{Demostración}~\cite{amable}:
    \begin{multline*}
        \mathbb{P}(N_1(t)=n,N_2(t)=m)\\
        =\mathbb{P}(N_1(t)=n,N_2(t)=m|\ N(t)=n+m)\cdot \mathbb{P}(N(t)=n+m)\\
        =\mathbb{P}(N_1(t)=n|\ N(t)=n+m)\cdot \mathbb{P}(N(t)=n+m)\\
        ={n+m\choose n}p^n(1-p)^m\cdot 
        \frac{(\lambda t)^{n+m}e^{-\lambda t}}{(n+m)!}\\
        =\frac{(n+m)!}{n!\ m!}p^n(1-p)^m\cdot 
        \frac{(\lambda t)^{n+m}e^{-\lambda t}}{(n+m)!}\\
        =\frac{(p\lambda\cdot t)^n e^{-p\lambda\cdot t}}{n!}
        \frac{[(1-p)\lambda\cdot t]^m e^{-(1-p)\lambda\cdot t}}{m!}\\
        = \mathbb{P}(N_1(t)=n)\cdot \mathbb{P}(N_2(t)=m)
    \end{multline*}
    con $N_1(t)\sim Poiss(p\lambda)$
    y $N_2(t)\sim Poiss((1-p)\lambda)$
\end{frame}


\begin{frame}{\secname: \subsecname}
    \textit{Ejemplo (cont.)}: ¿tenemos una red
    de Jackson en esta diapositiva?

    \begin{figure}
        \resizebox{.7\textwidth}{!}{%
            \usetikzlibrary{calc}
\begin{tikzpicture}


% the rectangular shape with vertical lines
\node[rectangle split, rectangle split parts=6,
draw, rectangle split horizontal,text height=1cm,text depth=0.5cm,inner ysep=0pt] (wa) {};
\fill[white] ([xshift=-\pgflinewidth,yshift=-\pgflinewidth]wa.north west) rectangle ([xshift=-15pt,yshift=\pgflinewidth]wa.south);

% the circle
\node[draw,circle,minimum size=1.5cm,anchor=west] (se) at (wa.east) {$\mu_1$};

% the arrows and labels
\draw[->] (se.east) -- +(20pt,0) node (seEnd) {};
\draw[<-] (wa.west) -- +(-20pt,0) node[left] {$\lambda_1$};




%%%%%%%%%%%%%%%%%%%
%%% Second queue
%%%%%%%%%%%%%%%%%%%

% the rectangular shape with vertical lines
\node[rectangle split, rectangle split parts=6,
    draw, rectangle split horizontal,text height=1cm,text depth=0.5cm,inner ysep=0pt,anchor=west] (wa2) at (seEnd.east) {};
\fill[white] ([xshift=-\pgflinewidth,yshift=-\pgflinewidth]wa2.north west) rectangle ([xshift=-15pt,yshift=\pgflinewidth]wa2.south);

% the circle
\node[draw,circle,minimum size=1.5cm,anchor=west] (se2) at (wa2.east) {$\mu_2$};

% the arrows and labels
\draw[->] (se2.east) -- +(40pt,0) node (seEnd2) {};



%%%%%%%%%%%%%%%%%%%
%%% Third queue
%%%%%%%%%%%%%%%%%%%

% the rectangular shape with vertical lines
\node[rectangle split, rectangle split parts=6,
    draw, rectangle split horizontal,text height=1cm,text depth=0.5cm,inner ysep=0pt,anchor=west] (wa3) at ($(wa2.south)-(2.5,1.5)$) {};
\fill[white]
    ([xshift=+\pgflinewidth,yshift=-\pgflinewidth]wa3.north east)
    rectangle ([xshift=-1pt,yshift=\pgflinewidth]wa3.south);
% the circle
\node[draw,circle,minimum size=1.5cm,anchor=east] (se3) at (wa3.west) {$\mu_3$};



%%%%%%%%%%%%%%%%%
%%% DRAW JOINS
%%%%%%%%%%%%%%%%%
\node[circle,inner sep=0pt,minimum size=3pt,draw=black,fill=black] (joinOut) at ($(se2.east)+(10pt,0)$) {};
\node[circle,inner sep=0pt,minimum size=3pt,draw=black,fill=black] (joinIn) at ($(wa.west)-(10pt,0)$) {};
\node[circle,inner sep=0pt,minimum size=3pt,draw=black,fill=black] (joinMid) at ($(wa2.west)-(15pt,0)$) {};
\node[circle,inner sep=0pt,minimum size=3pt,draw=black,fill=black] (joinMidDown) at ($(wa3.east)+(20pt,0)$) {};


% the exit
\draw[->] (joinOut) |- node[pos=0.3,anchor=west]
    {$p_{2,3}$} (wa3.east);

% the midfeedback
\draw[->] (joinMid) --
    ($(joinMid) - (0,1.1)$) -|
    node[pos=0.8,anchor=west] {$p_{1,3}$}
    (joinMidDown.north);

% the out feedback
\draw[->] (se3.west) -| (joinIn.south);



\end{tikzpicture}

        }
    \end{figure}

    Comprobemos las propiedades:
    \begin{enumerate}
        \item \cmark red abierta y
            ¿llegadas de Poisson?

            \begin{itemize}
                \item \cmark
                    cola $i=1$ con llegadas
                    $Poiss(\lambda_1)$;
                \item \cmark
                    cola $i=2$ con llegadas
                    $Poiss\left((1-p_{1,3})\lambda_1\right)$;
                \item \cmark
                    cola $i=3$ con llegadas
                    $Poiss\left(p_{1,3}\lambda_1+p_{2,3}\lambda_2\right)$;
            \end{itemize}
    \end{enumerate}
\end{frame}




\begin{frame}{\secname: \subsecname}

    \textit{Ejemplo (cont.)}: ¿tenemos una red
    de Jackson en esta diapositiva?

    \begin{figure}
        \resizebox{.7\textwidth}{!}{%
            \usetikzlibrary{calc}
\begin{tikzpicture}


% the rectangular shape with vertical lines
\node[rectangle split, rectangle split parts=6,
draw, rectangle split horizontal,text height=1cm,text depth=0.5cm,inner ysep=0pt] (wa) {};
\fill[white] ([xshift=-\pgflinewidth,yshift=-\pgflinewidth]wa.north west) rectangle ([xshift=-15pt,yshift=\pgflinewidth]wa.south);

% the circle
\node[draw,circle,minimum size=1.5cm,anchor=west] (se) at (wa.east) {$\mu_1$};

% the arrows and labels
\draw[->] (se.east) -- +(20pt,0) node (seEnd) {};
\draw[<-] (wa.west) -- +(-20pt,0) node[left] {$\lambda_1$};




%%%%%%%%%%%%%%%%%%%
%%% Second queue
%%%%%%%%%%%%%%%%%%%

% the rectangular shape with vertical lines
\node[rectangle split, rectangle split parts=6,
    draw, rectangle split horizontal,text height=1cm,text depth=0.5cm,inner ysep=0pt,anchor=west] (wa2) at (seEnd.east) {};
\fill[white] ([xshift=-\pgflinewidth,yshift=-\pgflinewidth]wa2.north west) rectangle ([xshift=-15pt,yshift=\pgflinewidth]wa2.south);

% the circle
\node[draw,circle,minimum size=1.5cm,anchor=west] (se2) at (wa2.east) {$\mu_2$};

% the arrows and labels
\draw[->] (se2.east) -- +(40pt,0) node (seEnd2) {};



%%%%%%%%%%%%%%%%%%%
%%% Third queue
%%%%%%%%%%%%%%%%%%%

% the rectangular shape with vertical lines
\node[rectangle split, rectangle split parts=6,
    draw, rectangle split horizontal,text height=1cm,text depth=0.5cm,inner ysep=0pt,anchor=west] (wa3) at ($(wa2.south)-(2.5,1.5)$) {};
\fill[white]
    ([xshift=+\pgflinewidth,yshift=-\pgflinewidth]wa3.north east)
    rectangle ([xshift=-1pt,yshift=\pgflinewidth]wa3.south);
% the circle
\node[draw,circle,minimum size=1.5cm,anchor=east] (se3) at (wa3.west) {$\mu_3$};



%%%%%%%%%%%%%%%%%
%%% DRAW JOINS
%%%%%%%%%%%%%%%%%
\node[circle,inner sep=0pt,minimum size=3pt,draw=black,fill=black] (joinOut) at ($(se2.east)+(10pt,0)$) {};
\node[circle,inner sep=0pt,minimum size=3pt,draw=black,fill=black] (joinIn) at ($(wa.west)-(10pt,0)$) {};
\node[circle,inner sep=0pt,minimum size=3pt,draw=black,fill=black] (joinMid) at ($(wa2.west)-(15pt,0)$) {};
\node[circle,inner sep=0pt,minimum size=3pt,draw=black,fill=black] (joinMidDown) at ($(wa3.east)+(20pt,0)$) {};


% the exit
\draw[->] (joinOut) |- node[pos=0.3,anchor=west]
    {$p_{2,3}$} (wa3.east);

% the midfeedback
\draw[->] (joinMid) --
    ($(joinMid) - (0,1.1)$) -|
    node[pos=0.8,anchor=west] {$p_{1,3}$}
    (joinMidDown.north);

% the out feedback
\draw[->] (se3.west) -| (joinIn.south);



\end{tikzpicture}

        }
    \end{figure}

    Comprobemos las propiedades:
    \begin{enumerate}
        \setcounter{enumi}{1}
        \item \cmark tiempos de servicio
            $Exp(\mu_i)$ y colas FIFO;
        \item \cmark vamos de $i$ a $j$
            con probabilidad $p_{i,j}$
            y salimos con probabilidad
            $1-\sum_j p_{i,j}$; y
        \item \cmark las cargas $\rho_i<1$
            si diseñamos la red para
            que $\lambda_i<\mu_i$
    \end{enumerate}
\end{frame}




\begin{frame}{\secname: \subsecname}
    \begin{thm}[Teorema de Jackson]
        Sea una red de Jackson con $m$ sistemas
        M/M/1, existe $\mathbb{P}(\mathbf{K})$
        y se expresa como:
        \begin{equation}
            \mathbb{P}(\mathbf{K})
            =\prod_i^m\mathbb{P}(N_i=K_i)\nonumber
            =\prod_i^m\pi_{K_i}
            =\prod_i^m(1-\rho_i)\rho_i^{K_i}
        \end{equation}
        con $\rho_i$ la carga de la cola $i$.
    \end{thm}
\end{frame}


\begin{frame}{\secname: \subsecname}
    \textit{Ejemplo}:
    quiero saber la probabilidad de tener
    $\mathbf{K}=(3,1,4)$ usuarios en
    equilibrio:
    \begin{align*}
        \mathbb{P}(\mathbf{K})&=
        \left(\mathbb{P}(N_1=3),
        \mathbb{P}(N_2=1),
        \mathbb{P}(N_3=4)\right)\\
        &=\prod_{i=1}^3 (1-\rho_i)\rho_i^{K_i}\\
        &=(1-\rho_1)\rho_1^3\cdot
        (1-\rho_2)\rho_2^1 \cdot
        (1-\rho_3)\rho_3^4
    \end{align*}

    \vfill

    \begin{figure}
        \resizebox{.7\textwidth}{!}{%
            \usetikzlibrary{calc}
\begin{tikzpicture}


% the rectangular shape with vertical lines
\node[rectangle split, rectangle split parts=6,
draw, rectangle split horizontal,
rectangle split part fill={white,white,white,white,HotPink3!20,HotPink3!20},
text height=1cm,text depth=0.5cm,inner ysep=0pt] (wa) {};
\fill[white] ([xshift=-\pgflinewidth,yshift=-\pgflinewidth]wa.north west) rectangle ([xshift=-15pt,yshift=\pgflinewidth]wa.south);

% the circle
\node[draw,circle,minimum size=1.5cm,anchor=west,fill=HotPink3!20] (se) at (wa.east) {$\mu_1$};

% the arrows and labels
\draw[->] (se.east) -- +(20pt,0) node (seEnd) {};
\draw[<-] (wa.west) -- +(-20pt,0) node[left] {};




%%%%%%%%%%%%%%%%%%%
%%% Second queue
%%%%%%%%%%%%%%%%%%%

% the rectangular shape with vertical lines
\node[rectangle split, rectangle split parts=6,
    draw, rectangle split horizontal,text height=1cm,text depth=0.5cm,inner ysep=0pt,anchor=west] (wa2) at (seEnd.east) {};
\fill[white] ([xshift=-\pgflinewidth,yshift=-\pgflinewidth]wa2.north west) rectangle ([xshift=-15pt,yshift=\pgflinewidth]wa2.south);

% the circle
\node[draw,circle,minimum size=1.5cm,anchor=west,fill=HotPink3!20] (se2) at (wa2.east) {$\mu_2$};

% the arrows and labels
\draw[->] (se2.east) -- +(40pt,0) node (seEnd2) {};



%%%%%%%%%%%%%%%%%%%
%%% Third queue
%%%%%%%%%%%%%%%%%%%

% the rectangular shape with vertical lines
\node[rectangle split, rectangle split parts=6,
    draw, rectangle split horizontal,
    rectangle split part fill={HotPink3!20,
    HotPink3!20,HotPink3!20,white},
    text height=1cm,text depth=0.5cm,inner ysep=0pt,anchor=west] (wa3) at ($(wa2.south)-(2.5,1.5)$) {};
\fill[white]
    ([xshift=+\pgflinewidth,yshift=-\pgflinewidth]wa3.north east)
    rectangle ([xshift=-1pt,yshift=\pgflinewidth]wa3.south);
% the circle
\node[draw,circle,minimum size=1.5cm,anchor=east,fill=HotPink3!20] (se3) at (wa3.west) {$\mu_3$};



%%%%%%%%%%%%%%%%%
%%% DRAW JOINS
%%%%%%%%%%%%%%%%%
\node[circle,inner sep=0pt,minimum size=3pt,draw=black,fill=black] (joinOut) at ($(se2.east)+(10pt,0)$) {};
\node[circle,inner sep=0pt,minimum size=3pt,draw=black,fill=black] (joinIn) at ($(wa.west)-(10pt,0)$) {};
\node[circle,inner sep=0pt,minimum size=3pt,draw=black,fill=black] (joinMid) at ($(wa2.west)-(15pt,0)$) {};
\node[circle,inner sep=0pt,minimum size=3pt,draw=black,fill=black] (joinMidDown) at ($(wa3.east)+(20pt,0)$) {};


% entrance lambda1
\node[anchor=south west] at (joinIn.east) {$\lambda_1$};


% the exit
\draw[->] (joinOut) |- node[pos=0.3,anchor=west]
    {$p_{2,3}$} (wa3.east);

% the midfeedback
\draw[->] (joinMid) --
    ($(joinMid) - (0,1.1)$) -|
    node[pos=0.8,anchor=west] {$p_{1,3}$}
    (joinMidDown.north);

% the out feedback
\draw[->] (se3.west) -| (joinIn.south);



\end{tikzpicture}

        }
    \end{figure}
\end{frame}


\begin{frame}{\secname: \subsecname}
    \textit{Ejemplo (cont.)}: miremos qué
    carga tiene cada cola\footnote{$\mu_2=\mu_3=\mu_1$.}\footnote{$\lambda_1=\lambda+\lambda_3$.}.

    \begin{itemize}
        \item carga cola $i=1$:
            $\rho_1=\tfrac{\lambda_1}{\mu_1}$;
        \item carga cola $i=2$:
            $\rho_2=\tfrac{(1-p_{1,3})\lambda_1}{\mu_2}=(1-p_{1,3})\rho_1$; y
        \item carga cola $i=3$:
            $\rho_3=\tfrac{p_{1,3}\lambda_1
            +p_{2,3}(1-p_{1,3})\lambda_1}{\mu_3}=[p_{1,3}+p_{2,3}(1-p_{1,3})]\rho_1$.
    \end{itemize}


    \begin{figure}
        \resizebox{.7\textwidth}{!}{%
            \usetikzlibrary{calc}
\begin{tikzpicture}


% the rectangular shape with vertical lines
\node[rectangle split, rectangle split parts=6,
draw, rectangle split horizontal,
rectangle split part fill={white,white,white,white,HotPink3!20,HotPink3!20},
text height=1cm,text depth=0.5cm,inner ysep=0pt] (wa) {};
\fill[white] ([xshift=-\pgflinewidth,yshift=-\pgflinewidth]wa.north west) rectangle ([xshift=-15pt,yshift=\pgflinewidth]wa.south);

% the circle
\node[draw,circle,minimum size=1.5cm,anchor=west,fill=HotPink3!20] (se) at (wa.east) {$\mu_1$};

% the arrows and labels
\draw[->] (se.east) -- +(20pt,0) node (seEnd) {};
\draw[<-] (wa.west) -- +(-20pt,0) node[left] {};




%%%%%%%%%%%%%%%%%%%
%%% Second queue
%%%%%%%%%%%%%%%%%%%

% the rectangular shape with vertical lines
\node[rectangle split, rectangle split parts=6,
    draw, rectangle split horizontal,text height=1cm,text depth=0.5cm,inner ysep=0pt,anchor=west] (wa2) at (seEnd.east) {};
\fill[white] ([xshift=-\pgflinewidth,yshift=-\pgflinewidth]wa2.north west) rectangle ([xshift=-15pt,yshift=\pgflinewidth]wa2.south);

% the circle
\node[draw,circle,minimum size=1.5cm,anchor=west,fill=HotPink3!20] (se2) at (wa2.east) {$\mu_2$};

% the arrows and labels
\draw[->] (se2.east) -- +(40pt,0) node (seEnd2) {};



%%%%%%%%%%%%%%%%%%%
%%% Third queue
%%%%%%%%%%%%%%%%%%%

% the rectangular shape with vertical lines
\node[rectangle split, rectangle split parts=6,
    draw, rectangle split horizontal,
    rectangle split part fill={HotPink3!20,
    HotPink3!20,HotPink3!20,white},
    text height=1cm,text depth=0.5cm,inner ysep=0pt,anchor=west] (wa3) at ($(wa2.south)-(2.5,1.5)$) {};
\fill[white]
    ([xshift=+\pgflinewidth,yshift=-\pgflinewidth]wa3.north east)
    rectangle ([xshift=-1pt,yshift=\pgflinewidth]wa3.south);
% the circle
\node[draw,circle,minimum size=1.5cm,anchor=east,fill=HotPink3!20] (se3) at (wa3.west) {$\mu_3$};



%%%%%%%%%%%%%%%%%
%%% DRAW JOINS
%%%%%%%%%%%%%%%%%
\node[circle,inner sep=0pt,minimum size=3pt,draw=black,fill=black] (joinOut) at ($(se2.east)+(10pt,0)$) {};
\node[circle,inner sep=0pt,minimum size=3pt,draw=black,fill=black] (joinIn) at ($(wa.west)-(10pt,0)$) {};
\node[circle,inner sep=0pt,minimum size=3pt,draw=black,fill=black] (joinMid) at ($(wa2.west)-(15pt,0)$) {};
\node[circle,inner sep=0pt,minimum size=3pt,draw=black,fill=black] (joinMidDown) at ($(wa3.east)+(20pt,0)$) {};


% entrance lambda1
\node[anchor=south west] at (joinIn.east) {$\lambda_1$};


% the exit
\draw[->] (joinOut) |- node[pos=0.3,anchor=west]
    {$p_{2,3}$} (wa3.east);

% the midfeedback
\draw[->] (joinMid) --
    ($(joinMid) - (0,1.1)$) -|
    node[pos=0.8,anchor=west] {$p_{1,3}$}
    (joinMidDown.north);

% the out feedback
\draw[->] (se3.west) -| (joinIn.south);



\end{tikzpicture}

        }
    \end{figure}
\end{frame}



\begin{frame}{\secname: \subsecname}
    \textit{Ejemplo (cont. II)}: tomando
    $p_{i,j}=\tfrac{1}{2},\ \forall (i,j)$ 
    tenemos $\rho_2=\tfrac{\rho_1}{2},
    \rho_3=\tfrac{3\rho_1}{4}$.
    Por tanto la probabilidad del estado
    $\mathbf{K}=(3,1,4)$ en la red de Jackson
    es

    \begin{align*}
        \mathbb{P}(\mathbf{K})&=
        (1-\rho_1)\rho_1^3\cdot
        (1-\rho_2)\rho_2^1 \cdot
        (1-\rho_3)\rho_3^4\\
        &= (1-\rho_1)\rho_1^3\cdot
        \left(1-\frac{\rho_1}{2}\right)\frac{\rho_1}{2}\cdot
        \left(1-\frac{3\rho_1}{4}\right)\left(
        \frac{3\rho_1}{4}\right)^4
    \end{align*}



    \begin{figure}
        \resizebox{.7\textwidth}{!}{%
            \usetikzlibrary{calc}
\begin{tikzpicture}


% the rectangular shape with vertical lines
\node[rectangle split, rectangle split parts=6,
draw, rectangle split horizontal,
rectangle split part fill={white,white,white,white,HotPink3!20,HotPink3!20},
text height=1cm,text depth=0.5cm,inner ysep=0pt] (wa) {};
\fill[white] ([xshift=-\pgflinewidth,yshift=-\pgflinewidth]wa.north west) rectangle ([xshift=-15pt,yshift=\pgflinewidth]wa.south);

% the circle
\node[draw,circle,minimum size=1.5cm,anchor=west,fill=HotPink3!20] (se) at (wa.east) {$\mu_1$};

% the arrows and labels
\draw[->] (se.east) -- +(20pt,0) node (seEnd) {};
\draw[<-] (wa.west) -- +(-20pt,0) node[left] {};




%%%%%%%%%%%%%%%%%%%
%%% Second queue
%%%%%%%%%%%%%%%%%%%

% the rectangular shape with vertical lines
\node[rectangle split, rectangle split parts=6,
    draw, rectangle split horizontal,text height=1cm,text depth=0.5cm,inner ysep=0pt,anchor=west] (wa2) at (seEnd.east) {};
\fill[white] ([xshift=-\pgflinewidth,yshift=-\pgflinewidth]wa2.north west) rectangle ([xshift=-15pt,yshift=\pgflinewidth]wa2.south);

% the circle
\node[draw,circle,minimum size=1.5cm,anchor=west,fill=HotPink3!20] (se2) at (wa2.east) {$\mu_2$};

% the arrows and labels
\draw[->] (se2.east) -- +(40pt,0) node (seEnd2) {};



%%%%%%%%%%%%%%%%%%%
%%% Third queue
%%%%%%%%%%%%%%%%%%%

% the rectangular shape with vertical lines
\node[rectangle split, rectangle split parts=6,
    draw, rectangle split horizontal,
    rectangle split part fill={HotPink3!20,
    HotPink3!20,HotPink3!20,white},
    text height=1cm,text depth=0.5cm,inner ysep=0pt,anchor=west] (wa3) at ($(wa2.south)-(2.5,1.5)$) {};
\fill[white]
    ([xshift=+\pgflinewidth,yshift=-\pgflinewidth]wa3.north east)
    rectangle ([xshift=-1pt,yshift=\pgflinewidth]wa3.south);
% the circle
\node[draw,circle,minimum size=1.5cm,anchor=east,fill=HotPink3!20] (se3) at (wa3.west) {$\mu_3$};



%%%%%%%%%%%%%%%%%
%%% DRAW JOINS
%%%%%%%%%%%%%%%%%
\node[circle,inner sep=0pt,minimum size=3pt,draw=black,fill=black] (joinOut) at ($(se2.east)+(10pt,0)$) {};
\node[circle,inner sep=0pt,minimum size=3pt,draw=black,fill=black] (joinIn) at ($(wa.west)-(10pt,0)$) {};
\node[circle,inner sep=0pt,minimum size=3pt,draw=black,fill=black] (joinMid) at ($(wa2.west)-(15pt,0)$) {};
\node[circle,inner sep=0pt,minimum size=3pt,draw=black,fill=black] (joinMidDown) at ($(wa3.east)+(20pt,0)$) {};


% entrance lambda1
\node[anchor=south west] at (joinIn.east) {$\lambda_1$};


% the exit
\draw[->] (joinOut) |- node[pos=0.3,anchor=west]
    {$p_{2,3}$} (wa3.east);

% the midfeedback
\draw[->] (joinMid) --
    ($(joinMid) - (0,1.1)$) -|
    node[pos=0.8,anchor=west] {$p_{1,3}$}
    (joinMidDown.north);

% the out feedback
\draw[->] (se3.west) -| (joinIn.south);



\end{tikzpicture}

        }
    \end{figure}
\end{frame}




\begin{frame}{\secname: \subsecname}
    \textit{Ejemplo (cont. III)}: en situación
    de equilibrio,
    probabilidad de estar en el estado
    $\mathbf{K}=(3,1,4)$ depende solo de
    $\rho_1$
    \begin{figure}
        \begin{tikzpicture}

\begin{axis}[every axis plot post/.append style={
  mark=none,domain=0:1,samples=40},
  axis x line*=bottom,
  axis y line*=left,
  %ymin=0,
  %ymax=8,
  %xmax=1,
  xlabel={$\rho_1$},
  ylabel={$\mathbb{P}(\mathbf{K})$},
  enlargelimits=upper,
  width=4in,
  height=2in]
  \addplot[ultra thick,smooth,color=HotPink2] {(1-x)*x^3*(1-x/2)*(x/2)*(1-3/4*x)*(3/4*x)^4} node [pos=.25,anchor=north west]
      {M/U/1};
\end{axis}


\end{tikzpicture}


    \end{figure}
\end{frame}




\subsection{Tiempo Medio de tránsito}
\begin{frame}{\secname: \subsecname}
    Los usuarios pueden seguir varias rutas
    (paths) $\{\mathcal{P}_1,\mathcal{P}_2,
    \ldots\}$, cada una
    de ellas con una probabilidad
    $\mathbb{P}\left( \mathcal{P}_i \right)$.
    \vfill
    \begin{figure}
        \usetikzlibrary{calc}
\begin{tikzpicture}


% the rectangular shape with vertical lines
\node[rectangle split, rectangle split parts=6,
draw, rectangle split horizontal,text height=1cm,text depth=0.5cm,inner ysep=0pt] (wa) {};
\fill[white] ([xshift=-\pgflinewidth,yshift=-\pgflinewidth]wa.north west) rectangle ([xshift=-15pt,yshift=\pgflinewidth]wa.south);

% the circle
\node[draw,circle,minimum size=1.5cm,anchor=west] (se) at (wa.east) {$\mu_1$};

% the arrows and labels
\draw[->] (se.east) -- +(20pt,0) node (seEnd) {};
\draw[<-] (wa.west) -- +(-20pt,0) node[left] {$\lambda_1$};




%%%%%%%%%%%%%%%%%%%
%%% Second queue
%%%%%%%%%%%%%%%%%%%

% the rectangular shape with vertical lines
\node[rectangle split, rectangle split parts=6,
    draw, rectangle split horizontal,text height=1cm,text depth=0.5cm,inner ysep=0pt,anchor=west] (wa2) at (seEnd.east) {};
\fill[white] ([xshift=-\pgflinewidth,yshift=-\pgflinewidth]wa2.north west) rectangle ([xshift=-15pt,yshift=\pgflinewidth]wa2.south);

% the circle
\node[draw,circle,minimum size=1.5cm,anchor=west] (se2) at (wa2.east) {$\mu_2$};

% the arrows and labels
\draw[->] (se2.east) -- +(40pt,0) node (seEnd2) {};



%%%%%%%%%%%%%%%%%%%
%%% Third queue
%%%%%%%%%%%%%%%%%%%

% the rectangular shape with vertical lines
\node[rectangle split, rectangle split parts=6,
    draw, rectangle split horizontal,text height=1cm,text depth=0.5cm,inner ysep=0pt,anchor=west] (wa3) at ($(wa2.south)-(2.5,1.5)$) {};
\fill[white]
    ([xshift=+\pgflinewidth,yshift=-\pgflinewidth]wa3.north east)
    rectangle ([xshift=-1pt,yshift=\pgflinewidth]wa3.south);
% the circle
\node[draw,circle,minimum size=1.5cm,anchor=east] (se3) at (wa3.west) {$\mu_3$};



%%%%%%%%%%%%%%%%%
%%% DRAW JOINS
%%%%%%%%%%%%%%%%%
\node[circle,inner sep=0pt,minimum size=3pt,draw=black,fill=black] (joinOut) at ($(se2.east)+(10pt,0)$) {};
\node[circle,inner sep=0pt,minimum size=3pt,draw=black,fill=black] (joinIn) at ($(wa.west)-(10pt,0)$) {};
\node[circle,inner sep=0pt,minimum size=3pt,draw=black,fill=black] (joinMid) at ($(wa2.west)-(15pt,0)$) {};
\node[circle,inner sep=0pt,minimum size=3pt,draw=black,fill=black] (joinMidDown) at ($(wa3.east)+(20pt,0)$) {};


% the exit
\draw[->] (joinOut) |- node[pos=0.3,anchor=west]
    {$p_{2,3}$} (wa3.east);

% the midfeedback
\draw[->] (joinMid) --
    ($(joinMid) - (0,1.1)$) -|
    node[pos=0.8,anchor=west] {$p_{1,3}$}
    (joinMidDown.north);

% the out feedback
\draw[->] (se3.west) -| (joinIn.south);


% PATHS
\draw[->,ultra thick,HotPink3]
    ($(wa.west)+(-1,.6)$) --
    ($(se2.east)+(1,.6)$) node[midway,above]
    {$\mathcal{P}_1$};


\end{tikzpicture}

    \end{figure}
\end{frame}


\begin{frame}{\secname: \subsecname}
    \begin{lema}[Tiempo Medio de Tránsito]
        El tiempo medio de tránsito por una
        ruta $\mathcal{P}_i$ de una red
        de Jackson en régimen estacionario es
        \begin{equation}
            \mathbb{E}_{\mathcal{P}_i}[T(t)]
            =\sum_{\iota\in \mathcal{P}_i}
            \mathbb{E}_{\iota}[T(t)]
            =\sum_{\iota\in \mathcal{P}_i}
            \frac{1}{\mu_\iota-\lambda_\iota}
        \end{equation}
    \end{lema}
\end{frame}


\begin{frame}{\secname: \subsecname}
    \textit{Ejemplo}: el tiempo medio de
    tránsito en la ruta $\mathcal{P}_1$ es

    \begin{equation}
        \mathbb{E}_{\mathcal{P}_1}[T(t)]=
        \frac{1}{\mu_1-\lambda_1}+
        \frac{1}{\mu_2-(1-p_{1,3})\lambda_1}
    \end{equation}

    La probabilidad de pasar por la ruta
    $\mathcal{P}_1$ es
    \begin{equation}
        \mathbb{P}(\mathcal{P}_1)=
        (1-p_{1,3})(1-p_{2,3})
    \end{equation}

    \begin{figure}
        \resizebox{.7\textwidth}{!}{%
            \usetikzlibrary{calc}
\begin{tikzpicture}


% the rectangular shape with vertical lines
\node[rectangle split, rectangle split parts=6,
draw, rectangle split horizontal,text height=1cm,text depth=0.5cm,inner ysep=0pt] (wa) {};
\fill[white] ([xshift=-\pgflinewidth,yshift=-\pgflinewidth]wa.north west) rectangle ([xshift=-15pt,yshift=\pgflinewidth]wa.south);

% the circle
\node[draw,circle,minimum size=1.5cm,anchor=west] (se) at (wa.east) {$\mu_1$};

% the arrows and labels
\draw[->] (se.east) -- +(20pt,0) node (seEnd) {};
\draw[<-] (wa.west) -- +(-20pt,0) node[left] {$\lambda_1$};




%%%%%%%%%%%%%%%%%%%
%%% Second queue
%%%%%%%%%%%%%%%%%%%

% the rectangular shape with vertical lines
\node[rectangle split, rectangle split parts=6,
    draw, rectangle split horizontal,text height=1cm,text depth=0.5cm,inner ysep=0pt,anchor=west] (wa2) at (seEnd.east) {};
\fill[white] ([xshift=-\pgflinewidth,yshift=-\pgflinewidth]wa2.north west) rectangle ([xshift=-15pt,yshift=\pgflinewidth]wa2.south);

% the circle
\node[draw,circle,minimum size=1.5cm,anchor=west] (se2) at (wa2.east) {$\mu_2$};

% the arrows and labels
\draw[->] (se2.east) -- +(40pt,0) node (seEnd2) {};



%%%%%%%%%%%%%%%%%%%
%%% Third queue
%%%%%%%%%%%%%%%%%%%

% the rectangular shape with vertical lines
\node[rectangle split, rectangle split parts=6,
    draw, rectangle split horizontal,text height=1cm,text depth=0.5cm,inner ysep=0pt,anchor=west] (wa3) at ($(wa2.south)-(2.5,1.5)$) {};
\fill[white]
    ([xshift=+\pgflinewidth,yshift=-\pgflinewidth]wa3.north east)
    rectangle ([xshift=-1pt,yshift=\pgflinewidth]wa3.south);
% the circle
\node[draw,circle,minimum size=1.5cm,anchor=east] (se3) at (wa3.west) {$\mu_3$};



%%%%%%%%%%%%%%%%%
%%% DRAW JOINS
%%%%%%%%%%%%%%%%%
\node[circle,inner sep=0pt,minimum size=3pt,draw=black,fill=black] (joinOut) at ($(se2.east)+(10pt,0)$) {};
\node[circle,inner sep=0pt,minimum size=3pt,draw=black,fill=black] (joinIn) at ($(wa.west)-(10pt,0)$) {};
\node[circle,inner sep=0pt,minimum size=3pt,draw=black,fill=black] (joinMid) at ($(wa2.west)-(15pt,0)$) {};
\node[circle,inner sep=0pt,minimum size=3pt,draw=black,fill=black] (joinMidDown) at ($(wa3.east)+(20pt,0)$) {};


% the exit
\draw[->] (joinOut) |- node[pos=0.3,anchor=west]
    {$p_{2,3}$} (wa3.east);

% the midfeedback
\draw[->] (joinMid) --
    ($(joinMid) - (0,1.1)$) -|
    node[pos=0.8,anchor=west] {$p_{1,3}$}
    (joinMidDown.north);

% the out feedback
\draw[->] (se3.west) -| (joinIn.south);


% PATHS
\draw[->,ultra thick,HotPink2]
    ($(wa.west)+(-1,.6)$) --
    ($(se2.east)+(1,.6)$) node[anchor=west]
    {$\mathcal{P}_1$};


\end{tikzpicture}
%
        }
    \end{figure}
\end{frame}


\begin{frame}[allowframebreaks]
        \frametitle{Referencias}
        \bibliographystyle{amsalpha}
        \bibliography{refs.bib}
\end{frame}





\end{document}
